\begin{proof}[Proof of \cref{obj:prop:three-level-witness}]
Let \(c\) be the observable capacity vector obtained from \(P_{\mathrm{obs}}^{\star}\) on the singleton covariate space \(\{x^{\star}\}\) and outcome levels \(\{0,1,2\}\).

\begin{enumerate}
\item Consider the latent allocation from \cref{obj:def:three-level-witness}, written with the canonical hidden-outcome completion
\[
\begin{array}{c|c|c|c}
(D_0,D_1) & (S_0,S_1) & (Y_0,Y_1) & \text{mass}\\ \hline
(0,1)&(1,1)&(0,0)&3/40\\
(0,1)&(1,1)&(1,1)&1/10\\
(0,1)&(1,1)&(2,2)&3/40\\
(0,1)&(0,1)&(0,0)&1/40\\
(0,1)&(0,1)&(0,1)&3/20\\
(0,1)&(0,1)&(0,2)&3/40\\
(0,1)&(0,0)&(0,0)&1/2 .
\end{array}
\]
The masses are nonnegative and sum to
\[
\frac{3}{40}+\frac{1}{10}+\frac{3}{40}
+\frac{1}{40}+\frac{3}{20}+\frac{3}{40}
+\frac12
=
1.
\]
Let \(Z\) be independent of these potential variables with
\(P(Z=0)=P(Z=1)=1/2\), and set \(D=D_Z\), \(S=S_D\), and \(Y=Y_D\) on selected units.  Every latent atom has \(D_0=0\), \(D_1=1\), and \(S_0\le S_1\).  The instrument propensity in the only cell is \(1/2\), so the overlap inequalities at \(\varepsilon_Z=1/4\) are
\[
0<\frac14<\frac12,
\qquad
\frac14\le \frac12\le \frac34 .
\]
The always-selected complier mass is
\[
\frac{3}{40}+\frac{1}{10}+\frac{3}{40}=\frac14>0.
\]
Thus the law satisfies the structural requirements collected in \cref{obj:def:structural-law-class}, with the single cell retained and weak selection direction \(S_0\le S_1\).

Under \(Z=0\), the selected outcome masses are \(3/40,1/10,3/40\), and the unselected mass is \(3/4\).  Multiplying by \(P(Z=0)=1/2\) gives the first four observed masses
\[
\frac{3}{80},\quad \frac{1}{20},\quad \frac{3}{80},\quad \frac38 .
\]
Under \(Z=1\), the selected outcome masses are
\[
\left(\frac{3}{40},\frac{1}{10},\frac{3}{40}\right)
+
\left(\frac{1}{40},\frac{3}{20},\frac{3}{40}\right)
=
\left(\frac{1}{10},\frac14,\frac{3}{20}\right),
\]
and the unselected mass is \(1/2\).  Multiplying by \(P(Z=1)=1/2\) gives the last four observed masses
\[
\frac{1}{20},\quad \frac18,\quad \frac{3}{40},\quad \frac14 .
\]
Hence this three-level slate system realizes \(P_{\mathrm{obs}}^{\star}\), has the displayed encoded full law, and satisfies \(D_0=0\), \(D_1=1\), and \(S_0\le S_1\) almost surely. % lean: wRealization

\item By \cref{obj:def:observable-capacities}, for \(i\in\{0,1,2\}\),
\[
\ell_i(x^{\star})
=
P(D=0,S=1,Y=i\mid Z=0,X=x^{\star})
-
P(D=0,S=1,Y=i\mid Z=1,X=x^{\star}).
\]
The second conditional probability is \(0\) for each \(i\), and the \(Z=0\) denominator is \(1/2\).  Therefore
\[
(\ell_0(x^{\star}),\ell_1(x^{\star}),\ell_2(x^{\star}))
=
\left(
\frac{(3/80)}{1/2},
\frac{(1/20)}{1/2},
\frac{(3/80)}{1/2}
\right)
=
\left(\frac{3}{40},\frac{1}{10},\frac{3}{40}\right).
\]
% lean: wLower

\item Similarly, for \(j\in\{0,1,2\}\),
\[
h_j(x^{\star})
=
P(D=1,S=1,Y=j\mid Z=1,X=x^{\star})
-
P(D=1,S=1,Y=j\mid Z=0,X=x^{\star}).
\]
The second conditional probability is \(0\) for each \(j\), and the \(Z=1\) denominator is \(1/2\).  Hence
\[
(h_0(x^{\star}),h_1(x^{\star}),h_2(x^{\star}))
=
\left(
\frac{(1/20)}{1/2},
\frac{(1/8)}{1/2},
\frac{(3/40)}{1/2}
\right)
=
\left(\frac{1}{10},\frac14,\frac{3}{20}\right).
\]
% lean: wUpper

\item Summing the displayed capacities gives
\[
q_0(x^{\star})=\frac{3}{40}+\frac{1}{10}+\frac{3}{40}=\frac14,
\qquad
q_1(x^{\star})=\frac{1}{10}+\frac14+\frac{3}{20}=\frac12.
\]
Thus
\[
m(x^{\star})=\min\{q_0(x^{\star}),q_1(x^{\star})\}=\frac14,
\qquad
\Delta q(x^{\star})=q_1(x^{\star})-q_0(x^{\star})=\frac14 .
\]
Using the threshold-cut formulas in \cref{obj:def:threshold-cuts}, the lower-cut correction is
\[
\min\{\Delta q(x^{\star}),0\}=0.
\]
The three prefix differences are
\[
\ell_{\le0}(x^{\star})-h_{\le0}(x^{\star})=-\frac{1}{40},
\]
\[
\ell_{\le1}(x^{\star})-h_{\le1}(x^{\star})=-\frac{7}{40},
\]
\[
\ell_{\le2}(x^{\star})-h_{\le2}(x^{\star})=-\frac14 .
\]
The maximum with \(0\) is therefore
\[
B_L(x^{\star})=0.
\]
For the upper cut,
\[
m(x^{\star})=\frac{10}{40},
\]
and the strict-prefix/tail sums are
\[
\ell_{<0}(x^{\star})+h_{>0}(x^{\star})=\frac25,
\]
\[
\ell_{<1}(x^{\star})+h_{>1}(x^{\star})=\frac{9}{40},
\]
\[
\ell_{<2}(x^{\star})+h_{>2}(x^{\star})=\frac{7}{40}.
\]
Taking the minimum gives
\[
B_U(x^{\star})
=
\min\left\{\frac{10}{40},\frac25,\frac{9}{40},\frac{7}{40}\right\}
=
\frac{7}{40}.
\]
% lean: wCapacityFacts

\item The capacity vector is valid because each displayed coordinate of \(\ell\) and \(h\) is nonnegative.  The singleton cell has observed weight
\[
p_{x^{\star}}=1\ge0,
\]
and the aggregate survivor-complier mass is
\[
M=\sum_x p_xm(x)=1\cdot \frac14=\frac14>0.
\]
Since the observed law is induced by the latent slate system constructed above, it lies in the observed-law domain.  By \cref{obj:def:identified-interval},
\[
\Theta_I(P_{\mathrm{obs}}^{\star})
=
\left[
\frac{p_{x^{\star}}B_L(x^{\star})}{M},
\frac{p_{x^{\star}}B_U(x^{\star})}{M}
\right]
=
\left[
\frac{0}{1/4},
\frac{(7/40)}{1/4}
\right]
=
\left[0,\frac{7}{10}\right].
\]
% lean: wInterval

\item The hypotheses of \cref{obj:thm:full-law-endpoint-attainment} hold for the realized three-level slate system: the structural restrictions and tie-safe survivor model hold by the first step, and \(M=1/4>0\) by the preceding step.  Applying that result to \(P_{\mathrm{obs}}^{\star}\), overlap \(1/4\), and the retained singleton cell gives full laws \(P^L\) and \(P^U\), with three-level slate systems, that induce \(P_{\mathrm{obs}}^{\star}\) and attain the two endpoint values
\[
0
\qquad\text{and}\qquad
\frac{7}{10}.
\]
% lean: wEndpointWitnesses
\end{enumerate}
Combining the six displayed conclusions gives all assertions of \cref{obj:prop:three-level-witness}. % lean: three_level_witness_sharp
\end{proof}
